\uuid{vQyt}
\titre{Formule de l'espérance}
\auteur{Maxime Nguyen}
\organisation{AMSCC}
\datecreate{2026-11-24}
\niveau{L2}
\module{Probabilité}
\chapitre{Probabilité continue}
\sousChapitre{Espérance et Fonction de répartition}
\difficulte{1}
\theme{espérance, mode}

\contenu{
\begin{question}
\qtype{onlyone}
Soit $X$ une variable de densité $f$ et de fonction de répartition $F$, admettant une espérance. Alors $\mathbb{E}(X)$ est égal à :
\begin{answers}
\bad{$\displaystyle\int_{-\infty}^{+\infty} f(x)\,\mathrm{d}x$}
\bad{l'abscisse du point où $f$ est maximale}
\bad{$\displaystyle\int_{-\infty}^{+\infty} x\,F(x)\,\mathrm{d}x$}
\good{$\displaystyle\int_{-\infty}^{+\infty} x\,f(x)\,\mathrm{d}x$}
\end{answers}
\begin{explanations}
A vaut toujours $1$. B est le mode, qui n'a pas de
raison d'être égal à l'espérance : pour la loi exponentielle de paramètre
$\lambda$, la densité est maximale en $0$ alors que l'espérance vaut
$\frac1\lambda$. C mélange densité et fonction de répartition.

Contre-sens visé : confondre espérance et « valeur la plus probable ».
\end{explanations}
\end{question}
}
